\section{Cumulative span}

The exact-word identities used in this section are proved in
Section~\ref{sec:app_wielandt_exact_word_support}. The stabilization and
spectral linear algebra behind the cumulative estimates are collected in
Section~\ref{sec:app_wielandt_cumulative_support}.

\begin{definition}[Word evaluation for a finite matrix family]
    \label{def:kraus_eval_word}
    \lean{Kraus.evalWord}
    \leanok
    Let $K=\{K^i\}_{i=0}^{d-1}$ be a finite family of $D\times D$ matrices.
    For a word $w=(i_1,\ldots,i_n)$, its \emph{word evaluation} is
    \begin{align}
        K^w=K^{i_1}\cdots K^{i_n},
        \notag
    \end{align}
    with the empty word evaluated as the identity matrix.
\end{definition}

\begin{definition}[Fixed-length vector span]\label{def:vector_spread_span}
    \lean{Kraus.vectorSpreadSpan}
    \leanok
    \uses{def:kraus_eval_word}
    The \emph{fixed-length vector span} at length~$n$ is
    \begin{align}
        H_n(K,\varphi)
         & =\spn_{\C}\{K^w\varphi:|w|=n\}\subseteq\C^D.
        \notag
    \end{align}
    This is $S_n(K)\ket{\varphi}$ in the notation
    of~\cite{Sanz2010Quantum}.
\end{definition}

\subsection{Paper primitivity and indices}

For the quantitative bounds, write $d'=\dim S_1(A)$ for the number of
linearly independent Kraus operators. The source definitions and
Proposition~3 of~\cite{Sanz2010Quantum} distinguish uniform vector spreading,
eventual exact-word spanning, and spectral strong irreducibility.

The next seven results use the following notation. For nonnegative integers
$d$ and $D$, let $K=\{K^i\}_{i=0}^{d-1}$ be a finite family of $D\times D$
matrices, and let $\E_K$ be its Kraus map.

\begin{theorem}[Iterated Kraus map on a rank-one matrix]
    \label{thm:kraus_iterate_rank_one_sum}
    \lean{Kraus.mapLM_pow_rankOne_eq_sum}
    \leanok
    \uses{def:kraus_map}
    For every choice of $d$, $D$, and $K$ as above, every $q\ge0$, and every
    $\varphi\in\C^D$,
    \begin{align}
        \E_K^q(\ketbra{\varphi}{\varphi})
        =\sum_{|w|=q}\ketbra{K^w\varphi}{K^w\varphi}.
        \notag
    \end{align}
\end{theorem}

\begin{proof}\leanok
    \uses{thm:kraus_map_iterate_word_expansion}
    Expand the Kraus representation of the $q$-fold iterate. Each summand is
    $K^w\ketbra{\varphi}{\varphi}(K^w)^\dagger
      =\ketbra{K^w\varphi}{K^w\varphi}$.
\end{proof}

\begin{theorem}[Positive rank-one image from vector spreading]
    \label{thm:kraus_rank_one_posdef_from_spreading}
    \lean{Kraus.mapLM_pow_rankOne_posDef_of_vectorSpreadSpan_eq_top}
    \leanok
    \uses{def:kraus_map, def:vector_spread_span}
    For every choice of $d$, $D$, and $K$ as above and every $q\ge0$, suppose
    $H_q(K,\varphi)=\C^D$ for every nonzero
    $\varphi\in\C^D$. Then
    $\E_K^q(\ketbra{\varphi}{\varphi})>0$ for every
    $\varphi\neq0$.
\end{theorem}

\begin{proof}\leanok
    \uses{thm:kraus_iterate_rank_one_sum, thm:finite_frame_posdef}
    The preceding theorem expresses the image as the frame operator of the
    spanning family $(K^w\varphi)_{|w|=q}$. A finite frame operator is positive
    definite exactly when its vectors span the ambient space.
\end{proof}

\begin{theorem}[Positive-semidefinite preservation by Kraus iterates]
    \label{thm:kraus_iterate_psd}
    \lean{Kraus.mapLM_pow_posSemidef}
    \leanok
    \uses{def:kraus_map}
    For every choice of $d$, $D$, and $K$ as above and every $n\ge0$, the
    iterate $\E_K^n$ maps positive-semidefinite
    matrices to positive-semidefinite matrices.
\end{theorem}

\begin{proof}\leanok
    \uses{thm:kraus_iterate_rank_one_sum}
    Decompose a positive-semidefinite matrix into a finite sum of rank-one
    positive matrices and apply Theorem~\ref{thm:kraus_iterate_rank_one_sum}
    to every summand.
\end{proof}

\begin{theorem}[Positivity improvement from vector spreading]
    \label{thm:kraus_positivity_improving_from_spreading}
    \lean{Kraus.mapLM_pow_positivityImproving_of_vectorSpreadSpan_eq_top}
    \leanok
    \uses{def:kraus_map, def:vector_spread_span}
    For every choice of $d$, $D$, and $K$ as above and every $q\ge0$, suppose
    $H_q(K,\varphi)=\C^D$ for every nonzero $\varphi$.
    Then $\E_K^q(\rho)>0$ for every nonzero positive-semidefinite matrix
    $\rho$.
\end{theorem}

\begin{proof}\leanok
    \uses{thm:kraus_rank_one_posdef_from_spreading, thm:kraus_iterate_psd}
    Write $\rho$ as a sum of rank-one positive matrices. At least one vector
    in this decomposition is nonzero, so its image is positive definite by
    Theorem~\ref{thm:kraus_rank_one_posdef_from_spreading}; all remaining
    images are positive semidefinite by Theorem~\ref{thm:kraus_iterate_psd}.
\end{proof}

\begin{theorem}[Definiteness of fixed points from vector spreading]
    \label{thm:kraus_fixed_point_posdef_from_spreading}
    \lean{Kraus.posDef_fixedPoint_of_vectorSpreadSpan_eq_top}
    \leanok
    \uses{def:kraus_map, def:vector_spread_span}
    For every choice of $d$, $D$, and $K$ as above and every $q\ge0$, suppose
    $H_q(K,\varphi)=\C^D$ for every nonzero $\varphi$. Then every nonzero
    positive-semidefinite fixed point of $\E_K$ is positive definite.
\end{theorem}

\begin{proof}\leanok
    \uses{thm:kraus_positivity_improving_from_spreading}
    A fixed point of $\E_K$ is fixed by $\E_K^q$, whose action on every
    nonzero positive-semidefinite matrix is positive definite.
\end{proof}

\begin{theorem}[Definiteness of fixed points of positive powers]
    \label{thm:kraus_power_fixed_point_posdef_from_spreading}
    \lean{Kraus.posDef_pow_fixedPoint_of_vectorSpreadSpan_eq_top}
    \leanok
    \uses{def:kraus_map, def:vector_spread_span}
    For every choice of $d$, $D$, and $K$ as above and every $q\ge0$, suppose
    $H_q(K,\varphi)=\C^D$ for every nonzero $\varphi$. For every $p>0$, each
    nonzero
    positive-semidefinite fixed point of $\E_K^p$ is positive definite.
\end{theorem}

\begin{proof}\leanok
    \uses{thm:kraus_positivity_improving_from_spreading, thm:kraus_iterate_psd}
    If $\E_K^p(\rho)=\rho$, then $\E_K^{pq}(\rho)=\rho$. Factor this iterate
    as $\E_K^q\E_K^{(p-1)q}$. The inner image is nonzero and positive
    semidefinite, so the outer application is positive definite.
\end{proof}

\begin{theorem}[Irreducibility from vector spreading]
    \label{thm:kraus_irreducible_from_spreading}
    \lean{Kraus.isIrreducibleMap_mapLM_of_vectorSpreadSpan_eq_top}
    \leanok
    \uses{def:kraus_map, def:vector_spread_span, def:irreducible_cp}
    For every choice of $d$, $D$, and $K$ as above and every $q\ge0$, if
    $H_q(K,\varphi)=\C^D$ for every nonzero $\varphi$, then the Kraus map
    $\E_K$ is irreducible.
\end{theorem}

\begin{proof}\leanok
    \uses{thm:kraus_invariant_compression_lower_zero}
    Let $P$ be an invariant orthogonal projection. Every word $K^w$ preserves
    $\operatorname{ran}P$. If $P\neq0$, choose a nonzero vector in this range.
    Its length-$q$ word orbit spans $\C^D$ by hypothesis and remains in
    $\operatorname{ran}P$, hence $P=\Id$.
\end{proof}

\begin{theorem}[Conjugate eigenmatrix of a positive map]
    \label{thm:positive-map-conj-eigenvector}
    \lean{Kraus.conjTranspose_eigenvector}
    \leanok
    \uses{def:positive_map}
    Let $E:\MN{D}\to\MN{D}$ be positive. If $E(X)=\mu X$, then
    $E(X^\dagger)=\overline\mu X^\dagger$.
\end{theorem}

\begin{proof}\leanok
    \uses{thm:positive_map_conjTranspose}
    Positivity implies that $E$ preserves adjoints. Taking the adjoint of
    $E(X)=\mu X$ gives the result.
\end{proof}

\begin{theorem}[Finite-order eigenmatrix is power-fixed]
    \label{thm:root-eigenvector-power-fixed}
    \lean{Kraus.pow_eigenvector_of_root}
    \leanok
    Let $E:\MN{D}\to\MN{D}$ be linear. If $E(X)=\mu X$ and $\mu^p=1$, then
    $E^p(X)=X$.
\end{theorem}

\begin{proof}\leanok
    \uses{thm:pow-on-eigenspace}
    Theorem~\ref{thm:pow-on-eigenspace} gives $E^p(X)=\mu^pX=X$.
\end{proof}

\begin{theorem}[Conjugate finite-order eigenmatrix is power-fixed]
    \label{thm:positive-map-conj-power-fixed}
    \lean{Kraus.pow_conjTranspose_eigenvector_of_root}
    \leanok
    \uses{def:positive_map}
    Let $E:\MN{D}\to\MN{D}$ be positive. If $E(X)=\mu X$ and $\mu^p=1$,
    then $E^p(X^\dagger)=X^\dagger$.
\end{theorem}

\begin{proof}\leanok
    \uses{thm:positive-map-conj-eigenvector,
        thm:root-eigenvector-power-fixed}
    The conjugate eigenvalue is $\overline\mu$, and
    $(\overline\mu)^p=\overline{\mu^p}=1$.
\end{proof}

\begin{theorem}[Trace of a non-fixed eigenmatrix]
    \label{thm:trace-preserving-eigenvector-trace-zero}
    \lean{Kraus.trace_eigenvector_eq_zero}
    \leanok
    \uses{def:trace_preserving}
    Let $E:\MN{D}\to\MN{D}$ preserve the trace. If $E(X)=\mu X$ and
    $\mu\neq1$, then $\tr(X)=0$.
\end{theorem}

\begin{proof}\leanok
    Trace preservation gives $\mu\tr(X)=\tr(E(X))=\tr(X)$. Since
    $\mu-1\neq0$, the trace vanishes.
\end{proof}

\begin{theorem}[Hermitian power-fixed point from a finite-order eigenmatrix]
    \label{thm:kraus-hermitian-power-fixed-point}
    \lean{Kraus.exists_hermitian_ne_zero_trace_zero_pow_fixedPoint}
    \leanok
    \uses{def:channel}
    Let $E$ be a channel, and suppose $E(X)=\mu X$ with $X\neq0$,
    $\mu\neq1$, and $\mu^p=1$. Then there is a nonzero Hermitian matrix $H$
    such that $\tr(H)=0$ and $E^p(H)=H$.
\end{theorem}

\begin{proof}\leanok
    \uses{thm:root-eigenvector-power-fixed,
        thm:positive-map-conj-power-fixed,
        thm:trace-preserving-eigenvector-trace-zero}
    At least one of $X+X^\dagger$ and $i(X^\dagger-X)$ is nonzero. Both are
    Hermitian, both have trace zero, and each is fixed by $E^p$.
\end{proof}

\begin{theorem}[Uniqueness of power-fixed positive matrices for a Kraus map]
    \label{thm:kraus-power-fixed-point-unique-from-spreading}
    \lean{Kraus.posSemidef_pow_fixedPoint_unique}
    \leanok
    \uses{def:kraus_map, def:vector_spread_span}
    Suppose $H_q(K,\varphi)=\C^D$ for every nonzero $\varphi$. If $p>0$ and
    $\rho,\sigma$ are nonzero positive-semidefinite fixed points of
    $\E_K^p$, then $\sigma=c\rho$ for some $c\in\C$.
\end{theorem}

\begin{proof}\leanok
    \uses{thm:kraus_power_fixed_point_posdef_from_spreading,
        thm:posdef_fixedpoint_proportional}
    Fixed-length spreading makes every nonzero positive-semidefinite fixed
    point of $\E_K^p$ positive definite. The critical-scalar argument then
    gives proportionality.
\end{proof}

\begin{theorem}[Powers of completely positive maps]
    \label{thm:cp-map-natural-power}
    \lean{IsCPMap.pow}
    \leanok
    \uses{def:cp_map}
    If $E$ is completely positive, then $E^p$ is completely positive for every
    $p\geq0$.
\end{theorem}

\begin{proof}\leanok
    Induct on $p$, using complete positivity of the identity map at $p=0$ and
    closure of completely positive maps under composition for the induction
    step.
\end{proof}

\begin{theorem}[Powers of channels]
    \label{thm:channel-natural-power}
    \lean{Kraus.isChannel_pow}
    \leanok
    \uses{def:channel}
    If $E$ is a channel, then $E^p$ is a channel for every $p\geq0$.
\end{theorem}

\begin{proof}\leanok
    \uses{thm:cp-map-natural-power}
    Complete positivity and trace preservation are both closed under
    composition and hold for the identity map.
\end{proof}

\begin{theorem}[Vanishing of Hermitian power-fixed points]
    \label{thm:kraus-hermitian-power-fixed-point-zero}
    \lean{Kraus.hermitian_pow_fixedPoint_eq_zero}
    \leanok
    \uses{def:kraus_map, def:vector_spread_span}
    Let $K$ be trace-preserving and suppose
    $H_q(K,\varphi)=\C^D$ for every nonzero $\varphi$. If $p>0$ and
    $H=H^\dagger$ satisfies $\tr(H)=0$ and $\E_K^p(H)=H$, then $H=0$.
\end{theorem}

\begin{proof}\leanok
    \uses{thm:channel-natural-power,
        thm:cesaro_fixedpoint,
        thm:kraus-power-fixed-point-unique-from-spreading,
        thm:positive_fixedpoint_decomp}
    Decompose $H$ as a difference of positive-semidefinite fixed points of the
    channel $\E_K^p$. Both are proportional to one nonzero fixed point. Their
    traces and $\tr(H)=0$ force equal coefficients, hence $H=0$.
\end{proof}

\begin{theorem}[Positive-definite fixed point from vector spreading]
    \label{thm:kraus_posdef_fixed_point_from_tp_spreading}
    \lean{Kraus.exists_posDef_fixedPoint_of_isTP_of_vectorSpreadSpan_eq_top}
    \leanok
    \uses{def:kraus_map, def:vector_spread_span}
    Let $D>0$, and let $K=\{K^i\}_{i=0}^{d-1}$ be a finite family of
    $D\times D$ matrices satisfying
    \begin{align}
        \sum_{i=0}^{d-1}(K^i)^\dagger K^i
         & =\Id.
        \notag
    \end{align}
    Suppose there is a common length $q\ge0$ such that
    $H_q(K,\varphi)=\C^D$ for every nonzero $\varphi\in\C^D$.
    Then there is a positive-definite matrix $\rho$ such that
    $\E_K(\rho)=\rho$.
\end{theorem}

\begin{proof}\leanok
    \uses{thm:cesaro_fixedpoint,
        thm:kraus_fixed_point_posdef_from_spreading}
    The normalization makes $\E_K$ a channel. The Ces\`aro fixed-point theorem,
    Theorem~\ref{thm:cesaro_fixedpoint}, gives a nonzero positive-semidefinite
    fixed point. The common fixed-length spreading hypothesis makes this fixed
    point positive definite by
    Theorem~\ref{thm:kraus_fixed_point_posdef_from_spreading}.
\end{proof}

\begin{theorem}[Primitivity from vector spreading]
    \label{thm:kraus_primitive_from_tp_spreading}
    \lean{Kraus.isPrimitive_mapLM_of_isTP_of_vectorSpreadSpan_eq_top}
    \leanok
    \uses{def:kraus_map, def:vector_spread_span, def:primitive_channel}
    Let $D>0$, and let $K=\{K^i\}_{i=0}^{d-1}$ be a finite family of
    $D\times D$ matrices satisfying
    \begin{align}
        \sum_{i=0}^{d-1}(K^i)^\dagger K^i
         & =\Id.
        \notag
    \end{align}
    Suppose there is a common length $q\ge0$ such that
    $H_q(K,\varphi)=\C^D$ for every nonzero $\varphi\in\C^D$.
    Then the Kraus map $\E_K$ is primitive: it has a nonzero fixed point and
    $1$ is its only peripheral eigenvalue.
\end{theorem}

\begin{proof}\leanok
    \uses{thm:cesaro_fixedpoint, thm:kraus_irreducible_from_spreading,
        thm:peripheral_roots_irreducible_channel,
        thm:kraus-hermitian-power-fixed-point,
        thm:kraus-hermitian-power-fixed-point-zero}
    The normalization makes $\E_K$ a channel, so
    Theorem~\ref{thm:cesaro_fixedpoint} gives a nonzero positive-semidefinite
    fixed point. This supplies the peripheral eigenvalue $1$.
    Theorem~\ref{thm:kraus_irreducible_from_spreading} gives irreducibility.
    Thus $\E_K$ is an irreducible channel. By
    Theorem~\ref{thm:peripheral_roots_irreducible_channel}, every peripheral
    eigenvalue $\mu$ is a root of unity; choose $p>0$ with
    $\mu^p=1$.

    Suppose $\mu\neq1$. Theorem~\ref{thm:kraus-hermitian-power-fixed-point}
    gives a nonzero Hermitian matrix $H$ with $\tr(H)=0$ and
    $\E_K^p(H)=H$. Theorem~\ref{thm:kraus-hermitian-power-fixed-point-zero}
    gives $H=0$, a contradiction. Thus $\mu=1$.
\end{proof}

\section{Nonzero trace product}

\begin{theorem}[Eigenvector extraction from a full cumulative span]
    \label{thm:eigenvector_from_full_cumulative_span}
    \lean{Kraus.exists_eigenvector_of_cumulativeSpan_eq_top}
    \leanok
    Let $T_N(K)$ be the span of the products of at most $N$ matrices from a
    finite family $K$. If $D>0$ and $T_N(K)=\MN{D}$, then there are a word
    $w$, a nonzero scalar $\mu$, and a nonzero vector $\varphi\in\C^D$ such
    that
    \begin{align}
        |w| & \le N, \notag\\
        K^w\varphi & = \mu\varphi. \notag
    \end{align}
\end{theorem}

\begin{proof}
    \leanok
    \uses{thm:eigenvector_from_trace}
    Since the identity lies in $T_N(K)$ and has nonzero trace, some word of
    length at most $N$ has nonzero trace. Such a matrix has a nonzero
    eigenvalue and a corresponding nonzero eigenvector by
    Theorem~\ref{thm:eigenvector_from_trace}.
\end{proof}

\section{Eigenvector spreading}

\begin{theorem}[Cumulative eigenvector spreading for finite families]
    \label{thm:kraus_cumulative_eigenvector_spreading}
    \lean{Kraus.eigenvector_spreading_of_cumulativeSpan_eq_top}
    \leanok
    \uses{def:kraus_eval_word}
    Let $K=\{K^i\}_{i=0}^{d-1}$ be a finite family of $D\times D$ matrices,
    where $D>0$, and let $\varphi\neq0$. If, for some $N$, words in $K$ of length at most $N$
    span $\MN{D}$, then the vectors $K^w\varphi$ with $|w|\leq D-1$ span
    $\C^D$.
\end{theorem}

\begin{proof}\leanok
    The cumulative vector spaces are monotone and have dimension at most $D$.
    Once two consecutive spaces agree, the common space is invariant under
    every $K^i$. Full cumulative matrix span then forces it to contain
    $\MN{D}\varphi=\C^D$. Since the initial space is nonzero, it reaches
    dimension $D$ by step $D-1$.
\end{proof}

\begin{lemma}[Eigenvector padding for finite families]
    \label{lem:kraus_eigenvector_padding}
    \lean{Kraus.cumulativeVectorSpan_eq_vectorSpreadSpan_of_eigenvector}
    \leanok
    \uses{def:kraus_eval_word, def:vector_spread_span}
    Suppose $K^{i_0}\varphi=\mu\varphi$ with $\mu\neq0$. For every $n$, the
    span of $K^w\varphi$ over words with $|w|\leq n$ equals
    $H_n(K,\varphi)$.
\end{lemma}

\begin{proof}\leanok
    Repeating the letter $i_0$ pads a word of length $m\leq n$ to length $n$.
    Its action on $\varphi$ gains the nonzero factor $\mu^{n-m}$, so every
    shorter-word vector belongs to the fixed-length span. The reverse inclusion
    is immediate.
\end{proof}

\begin{theorem}[Fixed-length spreading from eventual word-span fullness]
    \label{thm:eigenvector_spreading_eventually_full}
    \lean{Kraus.vectorSpreadSpan_eq_top_of_cumulativeVectorSpan_eq_top_of_eigenvector,
        Kraus.vectorSpreadSpan_eq_top_of_hasEventuallyFullWordSpan_of_eigenvector}
    \leanok
    \uses{def:kraus_eventually_full_word_span, def:vector_spread_span}
    Let $K=(K^i)_{i=0}^{d-1}$ be a finite family of $D\times D$ matrices, where
    $D>0$, and suppose that $S_n(K)=\MN{D}$ for every sufficiently large
    $n$. If $\varphi\neq0$, $\mu\neq0$, and
    $K^{i_0}\varphi=\mu\varphi$, then
    $H_{D-1}(K,\varphi)=\C^D$.
\end{theorem}

\begin{proof}\leanok
    \uses{thm:kraus_cumulative_eigenvector_spreading,
        lem:kraus_eigenvector_padding}
    Eventual fullness gives a level at which the cumulative matrix span is
    all of $\MN{D}$. The cumulative vector spaces therefore reach $\C^D$ by
    level $D-1$. The eigenvector relation then pads shorter words to length
    $D-1$ by nonzero powers of $\mu$.
\end{proof}

\section{Quantum Wielandt bounds}

Section~\ref{sec:app_wielandt_augmentation} proves the one-step augmentation
facts used in cases~(2) and~(3). The blocking identities used in the general
case are proved in Section~\ref{sec:app_wielandt_blocking_support}.

In this section we return to the notation of
\cite[Theorem~6.9]{Wolf2012Quantum}: $K_n$ is the complex linear span of the
Kraus operators of the $n$-th power of the channel, $k=\dim K_1$ is the Kraus
rank, and $q$ is the least integer such that $K_m=\MN{d}$ for every $m\geq q$.

\begin{definition}[Kraus rank and quantum Wielandt index]
    \label{def:wolf_kraus_rank_wielandt_index}
    \lean{Kraus.krausRank, Kraus.wielandtIndex}
    \leanok
    \uses{def:wolf_6_8_explicit_thresholds, def:kraus_exact_word_span}
    For a finite Kraus family on $\MN{d}$, define $k=\dim K_1$. When the
    family has eventually full word spans, define
    \begin{align}
        q&=\min\{q_0\in\mathbb N:K_m=\MN{d}\text{ for every }m\geq q_0\}.
        \notag
    \end{align}
    For a primitive quantum channel the set in the second definition is
    nonempty by Theorem~\ref{thm:wolf_6_8_kraus_equivalences}.
\end{definition}

\begin{lemma}[Least-threshold consequences]
    \label{lem:wolf_wielandt_index_minimal}
    \lean{Kraus.hasFullWordSpanFrom_wielandtIndex,
        Kraus.wielandtIndex_minimal,
        Kraus.wielandtIndex_le_of_wordSpan_eq_top_of_isTP}
    \leanok
    \uses{def:wolf_kraus_rank_wielandt_index}
    If the exact word spans are eventually full, then the least threshold
    $q$ is itself admissible. It is no larger than any other admissible
    threshold. In particular, for a trace-preserving Kraus family,
    $K_N=\MN{d}$ implies $q\leq N$.
\end{lemma}

\begin{proof}\leanok
    The set of admissible thresholds is nonempty by eventual fullness, so its
    infimum in $\mathbb N$ belongs to the set and is minimal. Trace preservation
    propagates the identity $K_N=\MN{d}$ to every $K_m$ with $m\geq N$.
\end{proof}

\begin{lemma}[A bounded word with nonzero trace]
    \label{lem:wolf_6_2_nonzero_trace_word}
    \lean{Kraus.wolf_lemma_6_2}
    \leanok
    \uses{def:wolf_kraus_rank_wielandt_index}
    Let $T:\MN{d}\to\MN{d}$ be a primitive quantum channel of Kraus rank
    $k$. Then there are an integer $n$ and a word $K^{(n)}\in K_n$ such that
    \begin{align}
        1\leq n\leq d^2-k+1,
        \qquad
        \tr K^{(n)}\neq0.
        \notag
    \end{align}
\end{lemma}

\begin{proof}\leanok
    \uses{thm:wolf_6_8_kraus_equivalences}
    Let $R_n$ be the span of the spaces $K_m$ with $1\leq m\leq n$.
    Starting from $\dim R_1=k$, strict growth continues until $R_n=\MN{d}$;
    otherwise equality at two consecutive levels would persist and contradict
    primitivity. Thus $R_{d^2-k+1}=\MN{d}$. Since the identity has nonzero
    trace, one of the spanning words occurring at a positive level has
    nonzero trace.
\end{proof}

\begin{lemma}[Spreading from a Kraus operator with nonzero eigenvalue]
    \label{lem:wolf_6_3_eigenvector_spreading}
    \lean{Kraus.rectSpan_nilpIndex_strict_growth_of_hasEventuallyFullWordSpan,
        Kraus.vecMulVec_eigenvector_mem_wordSpan_of_hasEventuallyFullWordSpan,
        Kraus.wolf_lemma_6_3}
    \leanok
    \uses{def:kraus_exact_word_span, def:vector_spread_span}
    Let $T:\MN{d}\to\MN{d}$ be a primitive quantum channel and suppose that
    one of its Kraus operators, denoted by $A$, satisfies
    $A\ket{\varphi}=\mu\ket{\varphi}$, where
    $\varphi\neq0$ and $\mu\neq0$. Then
    \begin{enumerate}
        \item $K_{d-1}\ket{\varphi}=\C^d$;
        \item if $A$ is not invertible, then
            $\ketbra{\varphi}{\psi}\in K_{d^2-d+1}$ for every
            $\psi\in\C^d$.
    \end{enumerate}
\end{lemma}

\begin{proof}\leanok
    \uses{thm:wolf_6_8_kraus_equivalences,
        thm:eigenvector_spreading_eventually_full}
    Part~(a) is Theorem~\ref{thm:eigenvector_spreading_eventually_full}.
    For part~(b), decompose $A$ into its nonzero and zero generalized
    eigenspaces. If $r$ is the largest zero Jordan-block size and
    $\widetilde d$ is the rank of $A^d$, strict growth of the rectangular
    spaces $A^rK_n$ reaches $A^r\MN{d}$ by level $d\widetilde d$.
    Hence $\ketbra{\varphi}{\psi}\in K_{d\widetilde d+r}$. Finally,
    $\widetilde d\leq d-r$ and $r\geq1$ give
    $d\widetilde d+r\leq d^2-d+1$.
\end{proof}

\begin{theorem}[Quantum Wielandt inequality]
    \label{thm:wolf_6_9_quantum_wielandt}
    \lean{Kraus.wolf_theorem_6_9,
        Kraus.wordSpan_eq_top_of_mem_wordSpan_one_of_isUnit,
        Kraus.wordSpan_eq_top_of_mem_wordSpan_one_of_nonzero_eigenvalue,
        Kraus.wielandtIndex_le_general,
        Kraus.wielandtIndex_le_of_mem_wordSpan_one_of_isUnit,
        Kraus.wielandtIndex_le_of_mem_wordSpan_one_of_nonzero_eigenvalue}
    \leanok
    \uses{def:wolf_kraus_rank_wielandt_index}
    Let $T:\MN{d}\to\MN{d}$ be a primitive quantum channel of Kraus rank
    $k$. Then the minimal threshold $q$ satisfies
    \begin{enumerate}
        \item in general, $q\leq(d^2-k+1)d^2$;
        \item if $K_1$ contains an invertible element, then
            $q\leq d^2-k+1$;
        \item if $K_1$ contains an element with a nonzero eigenvalue, then
            $q\leq d^2$.
    \end{enumerate}
\end{theorem}

\begin{proof}\leanok
    \uses{lem:wolf_6_2_nonzero_trace_word,
        lem:wolf_6_3_eigenvector_spreading,
        thm:kraus_word_span_upward_tp}
    For part~(2), adjoining an invertible element already in $K_1$ does not
    alter any $K_n$. Left multiplication by that element preserves linear
    independence, so the dimensions grow strictly from $k$ until they reach
    $d^2$. This gives $K_{d^2-k+1}=\MN{d}$.

    For part~(1), Lemma~\ref{lem:wolf_6_2_nonzero_trace_word} gives a word
    $K^{(n)}$ with nonzero trace and $n\leq d^2-k+1$. It has a nonzero
    eigenvalue. Regard the length-$n$ words as the Kraus operators of a
    blocked family and apply Lemma~\ref{lem:wolf_6_3_eigenvector_spreading}.
    The invertible branch uses part~(2); in the noninvertible branch the
    rank-one operators from part~(b), followed by the vector spreading in
    part~(a), give $K_{nd^2}=\MN{d}$. Thus
    $q\leq nd^2\leq(d^2-k+1)d^2$. Part~(3) is the same blocked argument with
    $n=1$.
\end{proof}
